\appendix

\section*{附录 A：给「怀疑自己是 ADHD」的自检线索}

这份清单不是诊断工具，它只做两件事：\term{帮你发现自己}，以及\term{帮你组织语言}（如果你打算去和专业人士聊）。逐条对照，不必每条都中——前面说过，伴随特征个体差异极大。

\subsection*{启动与任务}

\begin{checklist}
  \item 启动一件「没有兴趣、但有回报」的事，非常困难，常常拖到死线；
  \item 一旦没有外界任务压着，反而能自由地做自己的事、发展自己的想法；
  \item 对「必须做 X」这类命令有强烈的排异反应：犯困、想逃、刷手机；
  \item 任务「做完了没有」「做到哪一步了」很难自己追踪，需要别人提醒。
\end{checklist}

\subsection*{注意力}

\begin{checklist}
  \item 能连续数小时沉浸（超专注）在自己感兴趣的事上，但没法「按需专注」；
  \item 上课、开会、读长文时极容易走神，一走就拉不回来；
  \item 主动切换任务很痛苦，被动冒出念头跟着跑反而舒服；
  \item 经常「人醒着，但脑子像睡着了一样」：低唤醒、假性疲劳、午休后更累。
\end{checklist}

\subsection*{时间与组织}

\begin{checklist}
  \item 经常忘时间、错过约定、低估任务耗时；
  \item 生活层面功能崩塌：堆积的快递、没回的微信、乱到失控的房间；
  \item 需要靠死线、紧迫感才能启动；越没有压力越动不了。
\end{checklist}

\subsection*{情绪与冲动}

\begin{checklist}
  \item 对拒绝、批评、被评判高度敏感，容易过度反应或长期内耗；
  \item 容易上头也容易下头：三分钟热度、频繁开始又放弃；
  \item 用咖啡因、酒精、游戏、短视频、暴食等即时刺激自我调节；
  \item 心里有一本道德账本：长期用「我怎么这么懒」审判自己。
\end{checklist}

\subsection*{怎么使用这份清单}

\begin{enumerate}[leftmargin=2em,itemsep=2pt]
  \item 划掉那些和你无关的条目——这是正常的，伴随特征本来就因人而异；
  \item 找出 3--5 条对你生活\term{造成实际损害}的条目——那才是你的主线任务；
  \item 如果这些损害已经严重影响工作、关系或生活，去正规医院精神科或临床心理科聊聊。这不丢人，也和继续使用这份讲义不冲突；
  \item 无论是否确诊，第二、三部分的认知与方法都适用——它们依赖的是「你的大脑运行日志」，不是一张诊断书。
\end{enumerate}

\section*{附录 B：讲义黑话速查表}

\begin{center}
\small
\begin{tabular}{>{\raggedright\arraybackslash}p{3.2cm} p{12.0cm}}
\toprule
\textbf{黑话} & \textbf{含义} \\
\midrule
算法观 & 把 ADHD 理解为「广度便宜、深度昂贵；被动便宜、主动昂贵」的算法，而不是一处故障 \\
椭圆 & 脑内无意识念头的代称；超专注 = 某个念头占满注意力空间 \\
超专注 & 沉浸数小时、外界全屏蔽的状态；电源开关不在自己手上，但可引导 \\
归并效应 & 标签被广泛使用后，非必然的伴随特征也被贴上来的现象 \\
伴随特征 & 除执行功能缺陷外的附加特征，个体差异极大，甚至互相矛盾 \\
执行功能 & 维持线性集中 + 抑制其他念头（意志）；两者是一体的 \\
意志耗能 & 压住一个念头去干别的，可能比做那件事本身更费电 \\
恰当场 / 抽象轨道 & 把一堆事在内心合并成「一件事」，每天只在这一条宽轨道里被动切换 \\
程式性计划 & 计划「做 A」但不限定内容，随便翻，翻多了总会激活兴趣 \\
毛点计划 & 计划只是头脑里的锚点，靠环境刺激让「啪的一下」冒出来再行动 \\
框架欺骗 & 告诉自己「两小时干啥都行」，实际限制在 A 范围内，撤掉「必须」警报 \\
布局神经元 & 布置环境改变「猎物偏好」，降低启动门槛，以环境代意志 \\
冲动理性 & 冲动伴随新想法、神经生长与低耗能；顺着冲动调控成「网」而非「线」 \\
时间理性 & 默认「干 A 还是 B 只差切换成本」的普通理性模型，ADHD 上站不住脚 \\
动机悖论 & 社会规范压力推高皮质醇锁死启动；「人类要完蛋」这类念头反而消除压力、释放执行力 \\
隐形回血 & 有价值的创新性分心是回血；管理分心而非消除分心 \\
唤醒度三档 & 高→线性任务；中→翻找激发；低→躺下等想法冒出来再唤醒 \\
假性疲劳 & 缺的是定向能量而非追随性能量；「必须做 X」的念头本身在前额叶耗电 \\
动态自由度 & 想法上高自由度 + 操作上低自由度 \\
模糊分类 & 给学科建 4--5 个粗糙大类以限定广度搜索范围、节约算力 \\
慕容水战法 & 把任务当「或战或亡」的生死战，急性压力反而释放认知资源；不可常用 \\
蒙混过关法 & 外界线性任务先想「怎么蒙混过关」再迭代；心理负担大，少有人愿意用 \\
以广促深 & 广度算法在长周期里能解决深度问题，是贯穿全讲义的总心法 \\
\bottomrule
\end{tabular}
\end{center}

\clearpage
\section*{附录 C：出处对照}

本讲义的思想内核提炼自 B 站 UP 主「阿卡迪萨」（ADHD 野路子）的 24 期视频。以下为主要概念与原视频的对照，方便你回看原作：

\begin{small}
\begin{longtable}{>{\raggedright\arraybackslash}p{4.2cm} >{\raggedright\arraybackslash}p{11.0cm}}
\toprule
\textbf{视频（标题）} & \textbf{本讲义中的相关章节 / 概念} \\
\midrule
\endfirsthead
\toprule
\textbf{视频（标题）} & \textbf{本讲义中的相关章节 / 概念（续）} \\
\midrule
\endhead
\bottomrule
\endfoot
ADHD 是一套不同的算法 & 第一章算法观、椭圆竞争、以广促深；第十章 \\
ADHD 和 ASD 标签的「归并效应」 & 第一章核心特征与伴随特征、确诊 = 特征 + 损害 \\
ADHD 算是病吗？ & 第六章「正常人被极大高估」；第十章疾病观 \\
ADHD 学习的时候该听音乐吗？ & 第八章方法九（音乐策略） \\
ADHD 的执行功能问题及如何调控 & 第二章意志耗能；第八章恰当场、游戏管理 \\
ADHD 醒着但头脑已经「睡着」了 & 第二章低唤醒、唤醒度三档；第八章 \\
为什么午休后感觉更累？ & 第二章低唤醒与被动休息无效 \\
「时间理性」和「冲动理性」 & 第三章；第八章冲动引导成网 \\
游戏中无敌后的无聊现象 & 第三章兴趣引导；「激发场」 \\
「假性疲劳」与「动态自由度调节」 & 第二章；第八章方法八 \\
「程式性计划」但不限定内容 & 第八章方法二 \\
一个有效但没人愿意学的方法 & 第七章蒙混过关法；第九章 \\
应对主动任务切换 & 第二章主动 / 被动切换；第八章恰当场 \\
对「分心」的适当利用 & 第二章无聊优先；分心管理 \\
ADHD 的「动机悖论」 & 第三章动机悖论、一元论视角 \\
认知过载的异质性 & 第八章方法四（框架欺骗） \\
抑制可能比做事更耗能 & 第二章意志耗能；第八章 \\
隐形回血机制 & 第三章；第八章方法七 \\
古代人与现代人的压力；或战或亡 & 第八章方法十一（慕容水战法） \\
抗慢性压力的能力偏弱 & 第二章；第八章方法十一 \\
ADHD 可以做计划吗？ & 第八章方法三（毛点计划） \\
算法和抑制系统的补充（模糊分类） & 第八章方法十 \\
对 ADHD 好用的电脑软件 & 第八章方法十二（工具层） \\
在 B 站搞笑后发现自己好像是 ADHD & 算法观总述；第十章轴的程度、欧拉猜想 \\
\end{longtable}
\end{small}

\vspace{1em}

最后再重申一次，也是原视频一直坚持的态度：\term{以上大多是猜测，个体差异极大}。请把这份讲义当作一张地图，而不是一份处方——地图的作用是让你看清地形，路终究要你自己走，且只有你自己能走。

祝你，从读懂自己开始，慢慢活得顺一点。
